\subsection*{Chapter 1: Review of Stoichiometric Calculations \& Solution Chemistry}

\subsubsection*{Concentration of Solutions}

\textbf{Solubility (S):}
\begin{equation}
S = \frac{m}{q} \times 100
\end{equation}
\textit{S}: Solubility (g/100g solvent), \textit{m}: Mass of solute (g), \textit{q}: Mass of solvent (g).

\vspace{5mm}
\textbf{Titer (T):}
\begin{equation}
T_{g/mL} = \frac{m}{V} \quad \text{and} \quad T_{mg/mL} = \frac{m \times 1000}{V}
\end{equation}
\textit{T}: Titer of the solution (g/mL or mg/mL), \textit{m}: Mass of solute (g), \textit{V}: Volume of solution (mL).

\vspace{5mm}

\textbf{Percent Concentration (\%):}
\begin{align}
\%(w/w) &= \frac{m}{m+q} \times 100 \\
\%(w/v) &= \frac{m}{V} \times 100 \\
\%(v/v) &= \frac{V_x}{V} \times 100
\end{align}
\textit{m}: Mass of solute (g), \textit{q}: Mass of solvent (g), \textit{V}: Total volume of solution (mL), \textit{V\textsubscript{x}}: Volume of solute (mL).

\vspace{5mm}

\textbf{Trace Concentrations:}
\begin{align}
C(\text{ppm}) &= \frac{m}{m+q} \times 10^6 \\
C(\text{ppb}) &= \frac{m}{m+q} \times 10^9
\end{align}

\textbf{Molarity (C\textsubscript{M}) - Molar Concentration:}
\begin{equation}
C_M = \frac{m}{M \cdot V} \times 1000
\end{equation}
\textit{C\textsubscript{M}}: Molarity (mol/L), \textit{m}: Mass of solute (g), \textit{M}: Molar mass of solute (g/mol), \textit{V}: Volume of solution (mL).

\vspace{5mm}

\textbf{Molality (C\textsubscript{m}):}
\begin{equation}
C_m = \frac{m}{M \cdot q} \times 1000
\end{equation}
\textit{C\textsubscript{m}}: Molality (mol/kg), \textit{q}: Mass of solvent (g).

\vspace{5mm}
\textbf{Normality (C\textsubscript{N}) - Equivalent Concentration:}
\begin{equation}
C_N = \frac{m}{E \cdot V} \times 1000
\end{equation}
\textit{C\textsubscript{N}}: Normality (eq/L), \textit{E}: Equivalent weight of solute (g/eq).

\vspace{5mm}

\textbf{Equivalent Weight (EW):}
\begin{equation}
EW = \frac{FW}{n}
\end{equation}
\textit{EW}: Equivalent weight (g/eq), \textit{FW}: Formula weight (g/mol), \textit{n}: Number of reacting units (eq/mol).
\vspace{5mm}

\vspace{5mm}

\textbf{Box Method (Alligation) Formula and Principles:}

\textbf{1. Concentration of Higher Component ($a$) Principles:}
\begin{align}
\text{For Hydrated Salts: } & a = \frac{\text{M}_{\text{anhydrous}}}{\text{M}_{\text{hydrate}}} \times 100\% \\
\text{For Impure Solids: } & a = (100\% - \text{Impurities}\%)
\end{align}
\textit{a} is the $\text{w/w}\%$ concentration of the pure solute in the higher concentration component. This must be calculated if the starting material is not $100\%$ pure solute.

\vspace{5mm}

\textbf{2. Box Method Mass Ratio:}
\begin{equation}
\frac{m_a}{m_b} = \frac{c - b}{a - c}
\end{equation}
\textit{a}: Concentration of Component $a$ (higher concentration), \textit{b}: Concentration of Component $b$ (lower concentration), \textit{c}: Desired final concentration (intermediate), $m_a$: Mass of Component $b$ required, $m_b$: Mass of Component $a$ required.

\vspace{5mm}

\textbf{3. General Applicability:}
The Box Method applies to any mixture where two components of known concentration ($a$ and $b$) are mixed to achieve an intermediate concentration ($c$), provided the \textbf{solute and concentration units are consistent} across $a$, $b$, and $c$.

\vspace{5mm}
\textbf{Dilution:}
\begin{equation}
M_{\text{stock}} \times V_{\text{stock}} = M_{\text{diluted}} \times V_{\text{diluted}}
\end{equation}
\textit{M}: Molarity (mol/L), \textit{V}: Volume (mL or L).

\vspace{5mm}

\textbf{Interconversion of Concentration Units:}
\begin{align}
C_M &= \frac{10 \cdot C\% \cdot d}{M} \tag{1} \\
C_N &= \frac{10 \cdot C\% \cdot d}{E} \tag{2} \\
C_N &= C_M \cdot n \tag{3} \\
C_{g/L} &= C_M \cdot M_X \tag{4}
\end{align}

\noindent
\begin{itemize}
    \item $C_M$: \textbf{Molar concentration} (Molarity) in $\text{mol/L}$.
    \item $C_N$: \textbf{Normal concentration} (Normality) in $\text{eq/L}$.
    \item $C\%$: \textbf{Percent concentration} ($\text{w/w}$), mass of solute per $100\text{g}$ of solution.
    \item $d$: \textbf{Density} of the solution in $\text{g/mL}$.
    \item $M$: \textbf{Molar mass} of the solute in $\text{g/mol}$.
    \item $E$: \textbf{Equivalent mass} of the solute in $\text{g/eq}$.
    \item $n$: \textbf{Number of reacting units} (Equivalence factor) in $\text{eq/mol}$. Usually cation's charge in salts.
    \item $C_{g/L}$: \textbf{Concentration} in $\text{g/L}$.
    \item $M_X$: \textbf{Molar mass} of substance X in $\text{g/mol}$.
\end{itemize}
